\section{Implementation and Reproducibility}
\label{sec:impl}

\subsection{Code organisation}
The pipeline is organised under \texttt{src/} as
\texttt{config.py} (model IDs, prompts, paths), \texttt{data\_loader.py}
(deterministic subsampling under a fixed seed), \texttt{baseline.py}
(target-only autoregressive loop), \texttt{speculative.py} (vanilla
draft--verify loop following Leviathan
\emph{et al.}~\cite{Leviathan2023}), \texttt{drift\_speculative.py}
(DriftDiffuse drafter; Section~\ref{sec:driftdiffuse}),
\texttt{metrics.py}, \texttt{evaluate.py}, and \texttt{visualize.py}.
All measurement helpers (CUDA event-based latency, TTFT/TPOT decomposition,
seed initialisation, CSV logging) are isolated in \texttt{utils.py}.

\subsection{Models and quantisation}
The target is Qwen2.5-7B-Instruct~\cite{Qwen2025} loaded in 8-bit via
\texttt{bitsandbytes}~\cite{Dettmers2022Int8} with
\texttt{device\_map="auto"}. Drafts are Qwen2.5-0.5B and Qwen2.5-1.5B
in \texttt{fp16}. We use Hugging Face Transformers'
\texttt{generate(assistant\_model=...)} interface for vanilla speculative
decoding~\cite{HuggingFaceAssistant} to ensure that our acceptance numbers
are reproducible from the same library version reported in
\texttt{env\_spec.txt}.

\subsection{Sample manifests and seeds}
For each dataset we draw a fixed sample-ID manifest with \texttt{seed=42}
and persist it to \texttt{manifests/\{gsm8k,mmlu,xsum\}.json}; every
configuration is evaluated on exactly the same prompt set. For deterministic
runs we set \texttt{do\_sample=False}, \texttt{temperature=0},
\texttt{top\_p=1.0}; for stochastic runs we use \texttt{temperature=0.7},
\texttt{top\_p=0.9}. The two best configurations are re-run with seeds
$\{42,123,999\}$ to obtain the speedup standard deviation $\sigma_S$
reported in Section~\ref{sec:results}.

\subsection{Measurement protocol}
Latency is measured with \texttt{torch.cuda.Event} pairs around the
\texttt{generate} call, with one warmup sample per configuration to absorb
CUDA graph compilation. Per-sample we record
\texttt{sample\_id, task, latency, ttft, tpot, n\_tokens, output\_text}
plus the speculative-only fields
\texttt{n\_draft\_calls, n\_target\_calls, n\_accepted, n\_proposed},
from which $\bar{\alpha}$ and $\bar{B}_{\text{eff}}$ are derived directly
without distributional assumptions.

\subsection{Artefact release}
The complete code, manifests, raw \texttt{results/*.csv}, and the
DriftDiffuse training notebook are released under the Assignment~2
repository. The exact CUDA toolkit, driver, Python, and library versions
are pinned in \texttt{env\_spec.txt} alongside this manuscript.
